\documentclass[11pt]{article}
\usepackage[a4paper,margin=25mm]{geometry}
\usepackage{amsmath,amssymb}
\usepackage[T1]{fontenc}
\usepackage{lmodern}
\usepackage{hyperref}
\hypersetup{colorlinks=true,urlcolor=blue}
\setlength{\parindent}{0pt}
\setlength{\parskip}{7pt}
\title{The exact image size of every quadratic translate over a prime field}
\author{SucRunBug}
\date{1 October 2026}
\begin{document}
\maketitle
\textbf{Conjecture 00000000075 -- Verdict: FALSE.}

\textbf{Filed statement.} For $x\mapsto x^2+c$ on $\mathbb F_p$, outside exceptional $c$, the image size is asymptotic to $(1-\sqrt{\pi/8})p$. This is one clause of conjecture 00000000075.

\section*{Proof}
Fix any odd prime $p$ and any $c\in\mathbb F_p$. Equality of two squares in the field is equivalent to
\[x^2=y^2\quad\Longleftrightarrow\quad (x-y)(x+y)=0\quad\Longleftrightarrow\quad x=y\ \text{or}\ x=-y.\]
The square $0$ has one preimage. Every nonzero square has exactly two distinct preimages $x$ and $-x$, because the characteristic is odd. The $p-1$ nonzero elements therefore give $(p-1)/2$ distinct nonzero squares. The image of the square map has $(p+1)/2$ elements.

Translation by $c$ is a bijection of the field. Consequently, for \emph{every} $c$, with no exceptions,
\[|\operatorname{im}(x\mapsto x^2+c)|=\frac{p+1}{2},\qquad
\frac{|\operatorname{im}(x\mapsto x^2+c)|}{p}=\frac12+\frac1{2p}.\]
The normalized image size tends to $1/2$ as the odd primes tend to infinity. The conjectured coefficient is different:
\[1-\sqrt{\pi/8}<\frac12,\]
since $\pi>3>2$ implies $\pi/8>1/4$. Thus the normalized image size, which is always at least $1/2$, cannot tend to the claimed smaller constant. Equivalently, its ratio to the claimed main term cannot tend to $1$.

This refutes the image-size clause and hence the conjunction in the filed conjecture. Nothing here asserts a verdict on its separate claims about cyclic points or maximum tail depth. Excluding finitely or otherwise many values of $c$ cannot repair this clause, because the exact image-size formula holds for every $c$.

\section*{Lean formalization}
The Lean project proves the exact image-cardinality identity for every odd prime and every translate, proves the normalized lower bound $1/2$, and proves that no sequence of odd-prime instances can converge to the conjectured coefficient. It refutes the asymptotic assertion itself, not only a finite sample.

\textbf{Reproduction.} In the supplied \texttt{lean4/} project, run
\texttt{lake exe cache get}, then \texttt{lake build} and
\texttt{lake env lean Check.lean}. The pinned versions are Lean and Mathlib
\texttt{v4.33.0}. The supplied \texttt{reproduce.py} performs independent
finite sanity checks; these checks are not the proof of any infinite assertion.

\textbf{Source.} \url{https://github.com/The-Last-Math-Competition/The-Last-Math-Competition/blob/main/conjectures/00000000075.md}
\end{document}
